The energy consumption of logistics properties is relevant for a wide range of stakeholders. While users, owners, and facility managers are primarily interested in specific consumption data for reporting and calculation purposes, developers, financiers, brokers, energy suppliers, and benchmark service providers rely on data for property comparisons, accurate energy reporting, and consumption optimization.

 

In order to record the respective interests and data requirements in a structured manner, the project surveyed each stakeholder group about the implications of energy consumption data for their daily business, the associated pains and gains, and any data requirements they might have from other stakeholders. It was assumed that data would only be made available to external stakeholders if they could in turn cover their own data requirements.

 

Three pairs of stakeholders are listed here as examples, each of which has complementary data requirements and could complement each other well and with low barriers in a suitably designed data room (see Table~\ref{tbl:facroom} for the data room involving the facility manager).

\begin{table}
\caption{Overview of stakeholder data exchange pairs involving the facility manager.}
\label{tbl:facroom} 
\begin{tabular}{p{3.5cm} p{4cm} p{4cm} p{2cm}}
\textbf{Data flow} & \textbf{Data points} & \textbf{Granularity} & \textbf{Periodicity} \\
\hline

Facility Manager \textrightarrow{} Broker 
& Consumption data (for energy performance certificate) 
& Consumption types (electricity, water, gas) 
& annually \\

Broker \textrightarrow{} Facility Manager 
& User for new properties (as sales support) 
& Total property 
& once \\
\hline

Facility Manager \textrightarrow{} Tenant
& Consumption data 
& Consumption types (electricity, water, gas), as granular as possible 
& At least monthly, ideally every 15 minutes via dashboard \\

%Facility Manager \textrightarrow{} Tenant
& Blind loads\footnote{@@@link to Wikipedia? link to ontology?}
& Total property 
& once \\

%Facility Manager \textrightarrow{} Tenant
& Peak loads 
& Types of consumption: electricity (lighting, intralogistics, etc.) 
& As detailed as possible\footnote{@@@how is that a period? peak load needs to be measured precisely when it happened, precise timestamp, difference between when measured and when communicated; two timestamps in a way, measure timestamp is relevant, you can communicate every month } \\

Tenant \textrightarrow{} Facility Manager 
& Consumption data (for properties not under management by FM) 
& Total consumption per user 
& monthly \\

%Tenant \textrightarrow{} Facility Manager 
& Building specifics (area, age, shift regime, etc.) 
& Number, size, location 
& annually \\
\hline
\end{tabular}
\end{table}


\begin{table}
\caption{Overview of stakeholder data exchange pairs between investor and benchmark service provider.
The benchmark service provider takes the shared data points (building specifics and consumption data) as basis for the calculated consumption per $m^2$, which is in turn shared back.
}
\label{tbl:invroom}
\begin{tabular}{p{3.5cm} p{4cm} p{4cm} p{2cm}}
\textbf{Data flow} & \textbf{Data points} & \textbf{Granularity} & \textbf{Periodicity} \\
\hline
Investor \textrightarrow{} Benchmark service provider 
& Building specifics (area, age, shift regime, etc.) 
& Total property 
& once \\

%Investor 
%Benchmark service provider 
& Consumption data 
& Consumption types (electricity, water, gas) 
& annually \\

%\rowcolor{gray!20}
Benchmark service provider \textrightarrow{} Investor 

%& Min/max/median building specifics (area, age, shift regime\footnote{@@@shift regime could change over time, no? yes, might change, but is not very important; rather, the min/max/median building specifics is not shared back, but is needed to calculate the square meter aggregated value.}, etc.) 
%& Total property 
%& once \\

%Benchmark service provider \textrightarrow{} Investor 
& Min/max/median consumption per m\textsuperscript{2} per property 
& Consumption types (electricity, water, gas) 
& annually \\
\hline
\end{tabular}
\end{table}

This list can be supplemented by at least nine other identified pairs, each of which has expressed mutual data requirements in the context of energy consumption or property data for logistics real estate.

 

In any case, it is necessary to exchange data beyond the boundaries of one's own company, which, due to the data silos that have been common up to now, different data formats and granularities, and general reservations about sharing data, is only possible with a great deal of individual effort.

 

This is to be addressed and resolved through collaborative data sharing using the GRANERGIZE web app developed as part of the project.

\subsection{Data Room}

The term Data Room was chosen deliberately, as these Data Rooms are conceptualised as demarcated areas with clearly defined access conditions. Access to such a room is controlled by a gatekeeper. The Data Room therefore represents a protected and restricted area in which only approved participants are permitted to enter on the basis of a corresponding authorisation.

In the scenario considered here, the facility manager (FM) acts as the owner of his Data Room. Various tenants and brokers are granted access to this room in order to find out what data the FM provides at which level of granularity and with what periodicity. If the data offered by the FM meets their requirements, the actors involved can submit their own data offers in return and agree to exchange data if there is mutual consent. If, on the other hand, the available data does not meet the visitors’ expectations, they can leave the Data Room at any time without further interaction.
